\section{Relation to Existing Theory}\label{sec:literature}

\paragraph{Government standards and outsider incentives.}
\citet{Petropoulou2013} studies national
minimum-quality standards and the inability
in some circumstances to implement a
world-optimal standard by bilateral
adjustments without transfers.
\citet{Costinot2008} compares incomplete
institutional agreements on product
standards under national treatment
versus mutual recognition. Hence neither
governmental standard-setting failure nor
the existence of nonstandardized product
dimensions is an independent priority claim.

The closest governmental predecessors already establish that
regional standards may create incentives incompatible with wider
coordination. \citet{GandalShy2001} study recognition and
regional standards unions in three countries, with government
welfare combining domestic consumers and global profits of the
home firm. Their differentiated brands are fixed on a Salop
circle and compete in prices. Moreover, the absence of a
\emph{formal} union can coexist with worldwide mutual
recognition; their no-union regime is not the no-interface-link
$\SW$ benchmark in this paper.
\citet{TakaradaKawabata2025} model three-country
regional versus multilateral regulatory mutual recognition.
Under independently chosen national standards, regional
coordination can be unattractive, while outsiders can block
multilateral enlargement in other regimes. They also
consider exogenous product differentiation and standards
compliance costs. These are substantive prior results, not
merely similar institutional labels. Our additional margin
is a post-agreement Nash adjustment of \emph{residual} product
attributes that can reverse the country consent margins
under a fixed-depth, bilateral-first agenda; it is not the
discovery that an outsider may prefer to veto accession.

\paragraph{Endogenous differentiation and strategic trade policy.}
\citet{BranderSpencer2015} analyze endogenous
horizontal differentiation in international
Bertrand and Cournot oligopoly, including its
effect on patterns and gains from trade.
\citet{ZhouSpencerVertinsky2002} examine
strategic trade policy with endogenous export
quality, and \citet{JeanneretVerdier1996}
study trade-policy commitment and standards
in vertically differentiated markets.
Policy-dependent private product characteristics
and national exporter-profit effects are
therefore known mechanisms. The present
conditional response and incidence propositions
target the separate question of how the
\emph{same} staged government standardization
institution changes its country-specific
participation feasibility when firms
re-optimize nonstandardized horizontal attributes.

\paragraph{Firm compatibility, quality, and investment.}
The opposition between network benefits and product-market
competition is well established \citep{EconomidesFlyer1998,
EconomidesSkrzypacz2003}.
\citet{CabralSalant2014} establish that standards can reduce
private technology-development incentives. The published
article's proof numbering has not been independently
matched to the accessible earlier working-paper version;
we therefore do not attribute particular numbered
propositions to it.
\citet{BranderSpencer2022} also show that making
horizontal product differentiation endogenous can
reverse familiar Bertrand-versus-Cournot price and
profit rankings. Thus a disagreement between
fixed- and endogenous-differentiation forecasts
is not itself a new generic mechanism; their model
does not study three-country government accession
at inherited technical depth.
More closely, \citet{ChengHuang2025} examine endogenous
product quality and firm veto rights over compatibility in
network industries. Their published abstract reports
incompatibility under strong network effects when firms
hold veto power, together with positive welfare effects
of compatibility. Their accessible publisher appendix
discusses firm-level veto/no-veto subgames.
\citet{BuccellaFantiGori2024} likewise address
compatibility choice and quality differences in Cournot
network goods; a published corrigendum exists, so we
do not rely on an unchecked original theorem beyond
the available scope. Endogenous product characteristics,
private veto, and inefficient incompatibility are
therefore \emph{not} priority claims of the present paper.

\paragraph{Sequence and the conditional forecast.}
\citet{BuchholzHattori2024} examine technology choices
and coalition formation in a public-good game, with the
technology decision preceding coalition formation.
Here the government chooses interface compatibility
first; private firms subsequently optimize an independent,
costly horizontal-design margin; national governments anticipate the resulting domestic surplus
and exporters' rents when making their earlier consent decisions. The
particular result added to the predecessors is the
\emph{controlled forecast reversal} between endogenous
and frozen residual design, together with an inaccessible,
conditional Pareto gain, under the inherited-depth
A/B-first admission procedure.
Backward induction and uniqueness criteria draw on
standard methods \citep{RayVohra1999,Rosen1965}, not
new game-theoretic solution concepts.

This is a \emph{source-scoped} distinction, not a
comprehensive priority certificate. In particular,
the complete main-text propositions of
\citet{ChengHuang2025} were not available for exhaustive
proposition-level matching; this unresolved comparison
is disclosed rather than treated as evidence of novelty.
No result below establishes that all standardization
institutions, firm-design functions, or parameter vectors
exhibit the same failure.
